\chapter{Grupy rozwiązalne i rozszerzenia przez pierwiastniki}

\section{Szeregi pochodne grup}

\begin{definicja}
Komutator to $[g,h]=ghg^{-1}h^{-1}$. Podgrupa pochodna $G'=[G,G]$ jest podgrupą generowaną przez wszystkie komutatory. Definiujemy $G^{(0)}=G$, $G^{(i+1)}=(G^{(i)})'$. Grupa jest \emph{rozwiązalna}, jeżeli $G^{(r)}=\{1\}$ dla pewnego $r$.
\end{definicja}

Podgrupa $G'$ jest normalna, ponieważ sprzężenie komutatora jest komutatorem sprzężonych elementów. Dla $N\triangleleft G$ iloraz $G/N$ jest abelowy dokładnie wtedy, gdy $G'\subseteq N$: równość klas $ghN=hgN$ jest równoważna przynależności komutatora do $N$.

\begin{lemat}[Własności rozwiązalności]
Podgrupy i obrazy homomorficzne grup rozwiązalnych są rozwiązalne. Jeżeli $N\triangleleft G$, a $N$ i $G/N$ są rozwiązalne, to $G$ też jest rozwiązalna. Równoważnie, rozwiązalność oznacza istnienie szeregu
\[
G=G_0\triangleright G_1\triangleright\cdots\triangleright G_r=\{1\}
\]
o abelowych ilorazach $G_{i-1}/G_i$. Dla skończonej grupy można wybrać szereg, którego ilorazy są cykliczne rzędu będącego liczbą pierwszą.
\end{lemat}
\begin{proof}
Dla $H\le G$ indukcja daje $H^{(i)}\subseteq G^{(i)}$. Homomorfizm $\varphi$ wysyła komutatory na komutatory, więc $\varphi(G^{(i)})=\varphi(G)^{(i)}$. To dowodzi dwóch pierwszych własności. Jeśli $(G/N)^{(r)}=1$, to $G^{(r)}\subseteq N$. Jeśli ponadto $N^{(s)}=1$, to $G^{(r+s)}\subseteq N^{(s)}=1$.

Szereg pochodny grupy rozwiązalnej ma abelowe ilorazy. Odwrotnie, przy szeregu z abelowymi ilorazami indukcja daje $G^{(i)}\subseteq G_i$, bo $G_{i-1}'\subseteq G_i$. Zatem $G^{(r)}=1$.

Pozostaje udoskonalenie skończonego szeregu. Każda nietrywialna skończona grupa abelowa $A$ ma maksymalną właściwą podgrupę $B$. Grupa $A/B$ jest prosta i abelowa: każdy element różny od jedności generuje całość, a cykliczna grupa rzędu złożonego ma właściwą nietrywialną podgrupę. Stąd $A/B\simeq C_p$ dla pewnego pierwszego $p$. Powtarzając argument dla $B$ i biorąc przeciwobrazy podgrup w kolejnych abelowych ilorazach, otrzymujemy szereg z ilorazami cyklicznymi rzędu pierwszego.
\end{proof}

\begin{uwaga}
Jest to szereg subnormalny: $G_i$ ma być normalna w $G_{i-1}$, niekoniecznie w całym $G$. To wystarcza do zastosowania odpowiedniości Galois na każdym piętrze osobno.
\end{uwaga}

\begin{przyklad}[Dlaczego stopnie do czterech są szczególne]
Grupa $S_3$ ma szereg $S_3\triangleright A_3\triangleright\{1\}$ z ilorazami $C_2,C_3$. Dla $S_4$ używamy
\[
S_4\triangleright A_4\triangleright V_4\triangleright\{1\},\qquad
V_4=\{1,(12)(34),(13)(24),(14)(23)\}.
\]
Sprzężenie permutuje trzy podwójne transpozycje, więc $V_4$ jest normalna. Jej elementy różne od jedności mają rząd $2$, a iloczyn dwóch różnych takich elementów jest trzecim; stąd $V_4\simeq C_2\times C_2$. Ilorazy w szeregu mają rzędy $2,3,4$ i są abelowe. Zatem $S_4$, a także wszystkie podgrupy $S_n$ dla $n\le4$, są rozwiązalne.
\end{przyklad}

\section{Wieże pierwiastnikowe}

\begin{definicja}[Rozwiązalność przez pierwiastniki]
Wieża pierwiastnikowa to skończona wieża
\[
K=K_0\subseteq K_1\subseteq\cdots\subseteq K_r=M,\qquad
K_i=K_{i-1}(\alpha_i),\quad \alpha_i^{n_i}\in K_{i-1},\quad n_i\ge2.
\]
Wielomian jest rozwiązalny przez pierwiastniki nad $K$, jeśli jego ciało rozkładu zawiera się w końcowym ciele takiej wieży.
\end{definicja}

Dozwolone są działania arytmetyczne oraz wybór dowolnego zera $T^n-a$, także zespolonego. W szczególności wolno dołączać pierwiastki jedności. To szersza klasa operacji niż rzeczywiste pierwiastki kwadratowe występujące w konstrukcjach geometrycznych.

\section{Dwa lematy o dołączaniu pierwiastków}

Niech $\mu_n$ oznacza grupę zer $T^n-1$ w domknięciu algebraicznym. W charakterystyce zero ma ona $n$ elementów i jest cykliczna. Oto krótkie uzasadnienie używanej cykliczności. Każda skończona podgrupa $A$ grupy multiplikatywnej ciała jest abelowa. Jeśli $e$ jest najmniejszą wspólną wielokrotnością rzędów jej elementów, to $A$ zawiera element rzędu $e$: dla każdej potęgi pierwszej w $e$ wybieramy element, którego rząd zawiera tę potęgę, potęgowaniem usuwamy pozostałe czynniki i mnożymy otrzymane elementy o parami względnie pierwszych rzędach. Ich iloczyn ma rząd równy iloczynowi tych rzędów. Wszystkie elementy $A$ zerują $T^e-1$, więc $|A|\le e$. Element rzędu $e$ daje odwrotną nierówność, zatem generuje $A$. Generator $\mu_n$ oznaczamy przez $\zeta_n$.

\begin{lemat}
Rozszerzenie $K(\zeta_n)/K$ jest Galois z abelową grupą. Jeżeli $\mu_n\subseteq K$, $\alpha\ne0$ i $\alpha^n\in K$, to $K(\alpha)/K$ jest Galois, a jego grupa jest cykliczna rzędu dzielącego $n$.
\end{lemat}
\begin{proof}
Ciało $K(\zeta_n)$ jest ciałem rozkładu $T^n-1$. Automorfizm wysyła $\zeta_n$ na $\zeta_n^a$ z $\gcd(a,n)=1$, ponieważ zachowuje rząd multiplikatywny. Przypisanie automorfizmowi klasy $a$ zanurza grupę w abelowej grupie $(\Z/n\Z)^\times$.

Przy $\mu_n\subseteq K$ wszystkie zera $T^n-\alpha^n$ mają postać $\xi\alpha$, $\xi\in\mu_n$, i należą do $K(\alpha)$. Jest to zatem ciało rozkładu. Odwzorowanie $\sigma\mapsto\sigma(\alpha)/\alpha$ jest różnowartościowym homomorfizmem do $\mu_n$: elementy $\mu_n$ są ustalane przez $\sigma$, co zapewnia zachowanie mnożenia. Podgrupa grupy cyklicznej jest cykliczna.
\end{proof}

\begin{lemat}[Krok cykliczny rzędu pierwszego]
Niech $E/F$ będzie rozszerzeniem Galois z grupą cykliczną rzędu pierwszego $p$. Jeśli $\zeta_p\in F$, to istnieje $\alpha\in E$ takie, że $E=F(\alpha)$ oraz $\alpha^p\in F$.
\end{lemat}
\begin{proof}
Niech $\sigma$ generuje grupę. Z liniowej niezależności automorfizmów (lemat 3.1) operator
\[
R=\sum_{j=0}^{p-1}\zeta_p^{-j}\sigma^j
\]
nie jest zerowy. Wybierzmy $b\in E$ z $\alpha=R(b)\ne0$. Ponieważ $\sigma(\zeta_p)=\zeta_p$ i $\sigma^p=1$, przesunięcie indeksów daje $\sigma(\alpha)=\zeta_p\alpha$. Stąd $\alpha^p$ jest ustalane przez całą grupę, więc należy do $F$. Natomiast $\alpha\notin F$, bo $\zeta_p\ne1$. Stopień $[F(\alpha):F]>1$ dzieli pierwszy stopień $[E:F]=p$, więc jest równy $p$ i $F(\alpha)=E$.
\end{proof}

\section{Kryterium Galois: oba kierunki}

\begin{twierdzenie}[Kryterium pierwiastnikowe]
Niezerowy wielomian $f\in K[T]$ nad ciałem charakterystyki zero jest rozwiązalny przez pierwiastniki wtedy i tylko wtedy, gdy grupa jego ciała rozkładu jest rozwiązalna.
\end{twierdzenie}
\begin{proof}
Niech $L$ będzie ciałem rozkładu $f$.

\emph{Od pierwiastników do grupy.} Przypuśćmy, że $L\subseteq K_r$ dla wieży $K_i=K_{i-1}(\alpha_i)$, $\alpha_i^{n_i}=a_i\in K_{i-1}$. Usuwamy kroki z $a_i=0$, które nie powiększają ciała. Wybierzmy wspólną wielokrotność $N$ liczb $n_i$ i połóżmy $M_0=K(\zeta_N)$. Rozszerzenie $M_0/K$ jest Galois z abelową grupą.

Indukcyjnie zbudujemy skończone rozszerzenia Galois $M_i/K$ takie, że $K_i\subseteq M_i$ i $\Gal(M_i/K)$ jest rozwiązalna. Mając $M_{i-1}$, dołączamy wszystkie zera wielomianów
\[
T^{n_i}-\sigma(a_i),\qquad \sigma\in\Gal(M_{i-1}/K).
\]
Otrzymane ciało oznaczamy przez $M_i$. Zawiera $\alpha_i$ oraz $K_i$. Iloczyn wypisanych wielomianów ma współczynniki w $K$, bo grupa permutuje jego czynniki, a ciałem stałym jest $K$. Jeśli $g\in K[T]$ ma ciało rozkładu $M_{i-1}$, to $M_i$ jest ciałem rozkładu iloczynu $g$ i tego nowego wielomianu. Zatem $M_i/K$ jest rzeczywiście Galois.

Grupa $\Gal(M_i/M_{i-1})$ jest abelowa. Dla każdego nowego równania wybieramy jedno zero $\beta_\sigma$; pozostałe różnią się czynnikami z $\mu_{n_i}\subseteq M_{i-1}$. Odwzorowanie
\[
\tau\longmapsto\bigl(\tau(\beta_\sigma)/\beta_\sigma\bigr)_\sigma
\]
zanurza tę grupę w iloczynie grup $\mu_{n_i}$. Jest homomorfizmem, ponieważ wszystkie pierwiastki jedności są ustalane, i jest różnowartościowe, ponieważ wybrane zera generują rozszerzenie. Odpowiedniość Galois daje
\[
\Gal(M_i/K)/\Gal(M_i/M_{i-1})\simeq\Gal(M_{i-1}/K).
\]
Lemat 4.2 kończy indukcję. Na końcu $L\subseteq M_r$ i $L/K$ jest Galois, więc $\Gal(L/K)$ jest ilorazem rozwiązalnej grupy $\Gal(M_r/K)$. Jest zatem rozwiązalna.

\emph{Od grupy do pierwiastników.} Niech $G=\Gal(L/K)$ będzie rozwiązalna. Dla $G=1$ teza jest oczywista. W pozostałych przypadkach połóżmy $n=|G|$, $K'=K(\zeta_n)$ i $L'=LK'$, czyli najmniejsze ciało zawierające $L$ i $K'$. Rozszerzenie $L'/K'$ jest Galois, gdyż jest ciałem rozkładu $f$ nad $K'$. Ograniczenie automorfizmów do $L$ zanurza $\Gal(L'/K')$ w $G$: $L$ jest zachowywane jako ciało generowane przez zera $f$, a automorfizm ustalający $L$ i $K'$ ustala ich złożenie. Stąd grupa $G'=\Gal(L'/K')$ jest rozwiązalna.

Wybierzmy szereg subnormalny $G'=H_0\triangleright\cdots\triangleright H_r=1$ o ilorazach $H_{i-1}/H_i\simeq C_{p_i}$. Liczby pierwsze $p_i$ dzielą $|G'|$, a więc także $n$. Dla $E_i=(L')^{H_i}$ mamy
\[
K'=E_0\subseteq E_1\subseteq\cdots\subseteq E_r=L',
\qquad \Gal(E_i/E_{i-1})\simeq H_{i-1}/H_i\simeq C_{p_i}.
\]
Tutaj $E_i/E_{i-1}$ jest Galois, ponieważ $H_i\triangleleft H_{i-1}$, i stosujemy odpowiedniość Galois do rozszerzenia $L'/E_{i-1}$.
Każde $E_{i-1}$ zawiera $\zeta_{p_i}$ jako potęgę $\zeta_n$. Lemat 4.7 przedstawia więc każdy krok jako dołączenie pierwiastka $p_i$-tego stopnia. Pierwszy krok $K'/K$ też jest pierwiastnikowy, bo $\zeta_n^n=1$. Cała wieża jest pierwiastnikowa i jej koniec zawiera $L$, co dowodzi tezy.
\end{proof}

\begin{uwaga}[Dlaczego potrzebna była dodatkowa konstrukcja]
Z normalności kolejnych kroków wieży nie wynika normalność jej końca nad początkiem. W pierwszym kierunku dowodu dołączaliśmy dlatego wszystkie sprzężone pierwiastki, a nie tylko pierwiastki występujące w jednym wyrażeniu. Pominięcie tego kroku pozostawiałoby lukę w kryterium, na którym opiera się twierdzenie Abela--Ruffiniego.
\end{uwaga}

\begin{wniosek}
Każdy wielomian stopnia co najwyżej $4$ nad ciałem charakterystyki zero jest rozwiązalny przez pierwiastniki.
\end{wniosek}
\begin{proof}
Jego grupa Galois jest podgrupą $S_m$, gdzie $m\le4$ jest liczbą różnych zer. Wykazaliśmy rozwiązalność takich podgrup, więc stosujemy twierdzenie 4.8.
\end{proof}

Klasyczne wzory Cardana i Ferrariego są jawnymi wzorami dla stopni $3$ i $4$. Udowodnione kryterium wyjaśnia, dlaczego można je zbudować, nawet bez wypisywania ich rozbudowanej postaci.


